\section{Reproducible workflow and CI/CD}
\label{sec:cicd}
The repository uses continuous integration as a reproducibility mechanism, not as a substitute
for scientific interpretation. Figure~\ref{fig:cicd-workflow} shows the ordinary release path
and the two specialized evidence paths. The main workflow rebuilds the graph and static site,
while Fuseki and Silk run in separate workflows with their own runtime requirements.

\begin{figure}[htbp]\centering
\resizebox{\textwidth}{!}{\begin{tikzpicture}[
  node distance=8mm and 7mm,
  job/.style={draw=accent!80!black, rounded corners=3pt, fill=accent!5, align=center, font=\small, text width=3.0cm, minimum height=14mm},
  evidence/.style={draw=blue!70!black, rounded corners=3pt, fill=blue!7, align=center, font=\small, text width=3.0cm, minimum height=14mm},
  arr/.style={-{Stealth}, thick, draw=accent!80!black},
  branch/.style={-{Stealth}, dashed, thick, draw=black!55}
]
\node[job] (commit) at (0,0) {Git commit or pull request};
\node[job] (build) at (4.0,0) {Reconstruct from cache\\build static site\\compile report};
\node[job] (gates) at (8.0,0) {Quality gates\\reasoning + SHACL\\metrics + pytest};
\node[evidence] (artifact) at (12.0,0) {Release artifacts\\reports, RDF, PDF\\fingerprints};
\node[evidence] (pages) at (16.0,0) {GitHub Pages\\static publication};

\node[job] (fuseki) at (8.0,-2.7) {Fuseki/TDB2 job\\load, query, restart,\\read-only update checks};
\node[job] (silk) at (12.0,-2.7) {Silk job\\frozen candidate sample\\reference agreement};
\node[evidence] (review) at (16.0,-2.7) {Human review\\interpret limitations\\before accepting links};

\draw[arr] (commit)--(build)--(gates)--(artifact)--(pages);
\draw[branch] (gates)--(fuseki);
\draw[branch] (artifact)--(silk);
\draw[branch] (silk)--(review);
\node[font=\scriptsize, align=center] at (8,-1.15) {ordinary release path};
\node[font=\scriptsize, align=center] at (12,-3.75) {specialized evidence paths};
\end{tikzpicture}
}
\caption[CI/CD and evidence workflow]{CI/CD workflow. A commit enters the ordinary reconstruction
and quality gates; successful artifacts can be published to GitHub Pages. Fuseki/TDB2 and Silk
are separate evidence paths with different acceptance claims.}
\label{fig:cicd-workflow}
\end{figure}

\subsection{Release sequence}
The release sequence is intentionally ordered. Inputs are reconstructed first, RDF is then
materialized and checked, and only a release with matching fingerprints can be published or
loaded into Fuseki. A failure in a later gate does not turn a candidate graph into a public
graph. This design addresses a common failure mode in data projects: reporting a successful
transformation before validating the exact artifact that is served.

\begin{enumerate}
\item Reconstruct cached inputs and validate Silver records.
\item Generate asserted RDF, link metadata, ontology data, and selected inferred triples.
\item Run OWL diagnostics, SHACL, metrics, stale-release checks, and graph-isomorphism checks.
\item Build the static publication and the report from the accepted artifact.
\item Run the Fuseki/TDB2 acceptance job and, when requested, the isolated Silk experiment.
\end{enumerate}

\subsection{What CI proves}
The current suite demonstrates reproducible execution of the declared local gates. It does
not prove that every institution is represented, every identity link is correct, or that the
public endpoint can sustain an arbitrary workload. Those claims require independent annotation,
authoritative source comparison, and a load test against the selected production topology.
